% GENERATED FILE. Edit canonical lessons, then run npm run generate:books.
% canonical-source-hash: fnv1a64:e28cfa7d19629170
% canonical-lessons: PA-C293-mitti, PA-C293-zamin, PA-C293-paddhra, PA-C293-dhalan, PA-C293-tilla

\chapter{The Lie of the Land}
\label{ch:pa-a2-the-lie-of-the-land}
\hlchaptermodality{293}

\begin{chapteropening}
\textbf{By the end of this chapter:} \emph{I can describe the land around a village.}

Complete the last lesson of chapter 293: I can describe the land around a village.
\end{chapteropening}

\section[\texorpdfstring{\pa{ਮਿੱਟੀ} (miṭṭī)}{ਮਿੱਟੀ (miṭṭī)}]{\pa{ਮਿੱਟੀ} (miṭṭī) — soil; earth}

\label{lesson:PA-C293-mitti}



\begin{quote}
Before the new one: say the Punjabi for to dry up, then the Punjabi for to seep.
\end{quote}

\subsection*{You'll want to know: \pa{ਮਿੱਟੀ}}
\textbf{\pa{ਮਿੱਟੀ}} — \emph{miṭṭī} — \textquotedblleft{}soil; earth\textquotedblright{}.

\begin{quote}
\pa{ਸਾਡੇ} \pa{ਖੇਤ} \pa{ਦੀ} \pa{ਮਿੱਟੀ} \pa{ਨਰਮ} \pa{ਹੈ।}
\end{quote}

\begin{quote}
\emph{sāḍe khet dī miṭṭī naram hai.}
\end{quote}

\begin{quote}
The soil of our field is soft.
\end{quote}

A feminine noun.

\begin{cousinweb}
\textbf{\pa{ਮਿੱਟੀ}} (\emph{miṭṭī}) is inherited from Sanskrit \emph{mṛttikā}, \textquotedblleft{}earth, clay\textquotedblright{}; CDIAL lists Punjabi \emph{miṭṭī}, \textquotedblleft{}earth, soil, clay\textquotedblright{}, under it (CDIAL 10286).
\end{cousinweb}

\begin{grammarlens}[title={What you've built}]
One more word for this chapter.
\end{grammarlens}

\subsection*{Your turn}
\begin{itemize}
\raggedright
  \item \emph{Say it:} \emph{miṭṭī}
  \item \emph{Say it:} \emph{miṭṭī}, once more
  \item \emph{Say it:} \emph{sāḍe khet dī miṭṭī naram hai.} — The soil of our field is soft.
  \item \emph{Recall:} say the Punjabi for to dry up, then the Punjabi for to seep, then say \emph{miṭṭī} again
\end{itemize}

\begin{tcolorbox}[breakable,colback=teal!4,colframe=teal!35!black,title={Before you move on}]
What does \pa{ਮਿੱਟੀ} mean? (\textquotedblleft{}Soil; earth\textquotedblright{}.) Say it once more.
\end{tcolorbox}
\section[\texorpdfstring{\pa{ਜ਼ਮੀਨ} (zamīn)}{ਜ਼ਮੀਨ (zamīn)}]{\pa{ਜ਼ਮੀਨ} (zamīn) — land; the ground}

\label{lesson:PA-C293-zamin}



\begin{quote}
Before the new one: say the Punjabi for to seep, then the Punjabi for soil.
\end{quote}

\subsection*{You'll want to know: \pa{ਜ਼ਮੀਨ}}
\textbf{\pa{ਜ਼ਮੀਨ}} — \emph{zamīn} — \textquotedblleft{}land; the ground\textquotedblright{}.

\begin{quote}
\pa{ਮੇਰੇ} \pa{ਪਿਤਾ} \pa{ਜੀ} \pa{ਕੋਲ} \pa{ਬਹੁਤ} \pa{ਜ਼ਮੀਨ} \pa{ਹੈ।}
\end{quote}

\begin{quote}
\emph{mere pitā jī kol bahut zamīn hai.}
\end{quote}

\begin{quote}
My father has a lot of land.
\end{quote}

A feminine noun.

\begin{cousinweb}
\textbf{\pa{ਜ਼ਮੀਨ}} (\emph{zamīn}) is a loan, usually traced to Persian \emph{zamīn}, \textquotedblleft{}earth, land\textquotedblright{}.
\end{cousinweb}

\begin{grammarlens}[title={What you've built}]
One more word for this chapter.
\end{grammarlens}

\subsection*{Your turn}
\begin{itemize}
\raggedright
  \item \emph{Say it:} \emph{zamīn}
  \item \emph{Say it:} \emph{zamīn}, once more
  \item \emph{Say it:} \emph{mere pitā jī kol bahut zamīn hai.} — My father has a lot of land.
  \item \emph{Recall:} say the Punjabi for to seep, then the Punjabi for soil, then say \emph{zamīn} again
  \item \emph{Write it:} copy \textbf{\pa{ਜ਼ਮੀਨ}} once, checking the whole word
\end{itemize}

\begin{tcolorbox}[breakable,colback=teal!4,colframe=teal!35!black,title={Before you move on}]
What does \pa{ਜ਼ਮੀਨ} mean? (\textquotedblleft{}Land; the ground\textquotedblright{}.) Say it once more.
\end{tcolorbox}
\section[\texorpdfstring{\pa{ਪੱਧਰਾ} (paddhrā)}{ਪੱਧਰਾ (paddhrā)}]{\pa{ਪੱਧਰਾ} (paddhrā) — level; flat}

\label{lesson:PA-C293-paddhra}



\begin{quote}
Before the new one: say the Punjabi for soil, then the Punjabi for land.
\end{quote}

\subsection*{You'll want to know: \pa{ਪੱਧਰਾ}}
\textbf{\pa{ਪੱਧਰਾ}} — \emph{paddhrā} — \textquotedblleft{}level; flat\textquotedblright{}.

\begin{quote}
\pa{ਪਿੰਡ} \pa{ਦੇ} \pa{ਨੇੜੇ} \pa{ਜ਼ਮੀਨ} \pa{ਪੱਧਰੀ} \pa{ਹੈ।}
\end{quote}

\begin{quote}
\emph{piṇḍ de neṛe zamīn paddhrī hai.}
\end{quote}

\begin{quote}
Near the village the land is level.
\end{quote}

\begin{cousinweb}
\textbf{\pa{ਪੱਧਰਾ}} (\emph{paddhrā}) is inherited from a reconstructed \emph{*paddhara}, \textquotedblleft{}straight, level\textquotedblright{}; CDIAL lists Punjabi \emph{paddharā}, \textquotedblleft{}smooth and even\textquotedblright{}, under it (CDIAL 7767).
\end{cousinweb}

\begin{grammarlens}[title={What you've built}]
One more word for this chapter.
\end{grammarlens}

\subsection*{Your turn}
\begin{itemize}
\raggedright
  \item \emph{Say it:} \emph{paddhrā}
  \item \emph{Say it:} \emph{paddhrā}, once more
  \item \emph{Say it:} \emph{piṇḍ de neṛe zamīn paddhrī hai.} — Near the village the land is level.
  \item \emph{Recall:} say the Punjabi for soil, then the Punjabi for land, then say \emph{paddhrā} again
  \item \emph{Write it:} copy \textbf{\pa{ਪੱਧਰਾ}} once, checking the whole word
\end{itemize}

\begin{tcolorbox}[breakable,colback=teal!4,colframe=teal!35!black,title={Before you move on}]
What does \pa{ਪੱਧਰਾ} mean? (\textquotedblleft{}Level; flat\textquotedblright{}.) Say it once more.
\end{tcolorbox}
\section[\texorpdfstring{\pa{ਢਲਾਣ} (ḍhalāṇ)}{ਢਲਾਣ (ḍhalāṇ)}]{\pa{ਢਲਾਣ} (ḍhalāṇ) — a slope}

\label{lesson:PA-C293-dhalan}



\begin{quote}
Before the new one: say the Punjabi for land, then the Punjabi for level.
\end{quote}

\subsection*{You'll want to know: \pa{ਢਲਾਣ}}
\textbf{\pa{ਢਲਾਣ}} — \emph{ḍhalāṇ} — \textquotedblleft{}a slope\textquotedblright{}.

\begin{quote}
\pa{ਪਹਾੜ} \pa{ਦੀ} \pa{ਢਲਾਣ} \pa{ਉੱਤੇ} \pa{ਰੁੱਖ} \pa{ਹਨ।}
\end{quote}

\begin{quote}
\emph{pahāṛ dī ḍhalāṇ utte rukkh han.}
\end{quote}

\begin{quote}
There are trees on the slope of the mountain.
\end{quote}

A feminine noun.

\begin{cousinweb}
\textbf{\pa{ਢਲਾਣ}} (\emph{ḍhalāṇ}) is built on \emph{ḍhalṇā}, \textquotedblleft{}to fall, be poured out\textquotedblright{}, from a reconstructed \emph{*ḍhalati}, \textquotedblleft{}bends over, falls\textquotedblright{}; CDIAL lists Punjabi \emph{ḍhalṇā}, \textquotedblleft{}to be poured out, fall\textquotedblright{}, under it (CDIAL 5581).
\end{cousinweb}

\begin{grammarlens}[title={What you've built}]
One more word for this chapter.
\end{grammarlens}

\subsection*{Your turn}
\begin{itemize}
\raggedright
  \item \emph{Say it:} \emph{ḍhalāṇ}
  \item \emph{Say it:} \emph{ḍhalāṇ}, once more
  \item \emph{Say it:} \emph{pahāṛ dī ḍhalāṇ utte rukkh han.} — There are trees on the slope of the mountain.
  \item \emph{Recall:} say the Punjabi for land, then the Punjabi for level, then say \emph{ḍhalāṇ} again
  \item \emph{Write it:} copy \textbf{\pa{ਢਲਾਣ}} once, checking the whole word
\end{itemize}

\begin{tcolorbox}[breakable,colback=teal!4,colframe=teal!35!black,title={Before you move on}]
What does \pa{ਢਲਾਣ} mean? (\textquotedblleft{}A slope\textquotedblright{}.) Say it once more.
\end{tcolorbox}
\section[\texorpdfstring{\pa{ਟਿੱਲਾ} (ṭillā)}{ਟਿੱਲਾ (ṭillā)}]{\pa{ਟਿੱਲਾ} (ṭillā) — a hillock; a mound}

\label{lesson:PA-C293-tilla}



\begin{quote}
Before the new one: say the Punjabi for level, then the Punjabi for a slope.
\end{quote}

\subsection*{You'll want to know: \pa{ਟਿੱਲਾ}}
\textbf{\pa{ਟਿੱਲਾ}} — \emph{ṭillā} — \textquotedblleft{}a hillock; a mound\textquotedblright{}.

\begin{quote}
\pa{ਪਿੰਡ} \pa{ਦੇ} \pa{ਪਿੱਛੇ} \pa{ਇੱਕ} \pa{ਉੱਚਾ} \pa{ਟਿੱਲਾ} \pa{ਹੈ।}
\end{quote}

\begin{quote}
\emph{piṇḍ de pichchhe ikk uchchā ṭillā hai.}
\end{quote}

\begin{quote}
There is a high mound behind the village.
\end{quote}

A masculine noun.

\begin{cousinweb}
\textbf{\pa{ਟਿੱਲਾ}} (\emph{ṭillā}) is inherited from a reconstructed \emph{*ṭalla}, \textquotedblleft{}heap\textquotedblright{}; CDIAL lists Punjabi \emph{ṭillā}, \textquotedblleft{}hillock, mound\textquotedblright{}, under it (CDIAL 5451).
\end{cousinweb}

\begin{grammarlens}[title={What you've built}]
One more word for this chapter.
\end{grammarlens}

\subsection*{Your turn}
\begin{itemize}
\raggedright
  \item \emph{Say it:} \emph{ṭillā}
  \item \emph{Say it:} \emph{ṭillā}, once more
  \item \emph{Say it:} \emph{piṇḍ de pichchhe ikk uchchā ṭillā hai.} — There is a high mound behind the village.
  \item \emph{Recall:} say the Punjabi for level, then the Punjabi for a slope, then say \emph{ṭillā} again
  \item \emph{Write it:} copy \textbf{\pa{ਟਿੱਲਾ}} once, checking the whole word
\end{itemize}

\begin{tcolorbox}[breakable,colback=teal!4,colframe=teal!35!black,title={Before you move on}]
What does \pa{ਟਿੱਲਾ} mean? (\textquotedblleft{}A hillock; a mound\textquotedblright{}.) Say it once more.
\end{tcolorbox}
